\thispagestyle{plain}
\chapter*{Resumen}

La selección óptima de carteras según el modelo media-varianza de Markowitz
depende de la estimación de dos parámetros: el vector de rentabilidades
esperadas $\boldsymbol{\mu}$ y la matriz de covarianza $\boldsymbol{\Sigma}$.
Este trabajo formula el modelo como un problema cuadrático convexo (QP), lo
resuelve de forma exacta a través de sus condiciones de Karush-Kuhn-Tucker (KKT)
y evalúa hasta qué punto varias técnicas estadísticas y de \textit{Machine
Learning} mejoran la estimación de esos parámetros. Sobre una muestra de 30
acciones del S\&P~500 (2010--2024), se comparan tres estimadores de la covarianza
(muestral, \textit{shrinkage} de Ledoit-Wolf y reconstrucción por componentes
principales, PCA) y tres modelos de predicción de $\boldsymbol{\mu}$ (Ridge,
Lasso y Random Forest), validados mediante \textit{backtesting} walk-forward sin
\textit{look-ahead} frente a los \textit{baselines} $1/N$ (equiponderada) y el
índice de mercado (SPY). El \textit{shrinkage} y el PCA mejoran de forma modesta
la cartera muestral (Sharpe $0{,}68$ y $0{,}69$ frente a $0{,}64$). Entrenados de
forma agrupada, los modelos de \textit{Machine Learning} resultan competitivos:
Lasso (Sharpe $0{,}75$) llega a igualar a la estrategia $1/N$ ($0{,}76$). Aun
así, ninguna cartera optimizada supera al mercado (Sharpe $0{,}90$) y, sobre
todo, ninguna diferencia de Sharpe es estadísticamente significativa. La
principal contribución no es un modelo que bata al mercado, sino una
demostración rigurosa de por qué el error de estimación domina el problema,
hasta hacer indistinguibles las estrategias, y de qué técnicas ayudan a
mitigarlo.

\keywordses{Optimización de carteras, Modelo de Markowitz, Programación cuadrática convexa, Shrinkage de Ledoit-Wolf, Machine Learning, Backtesting walk-forward.}

\MediaOptionLogicBlank

\pdfbookmark[1]{Abstract}{abstract}
\chapter*{Abstract}
Optimal portfolio selection under Markowitz's mean--variance model depends on
estimating two parameters: the vector of expected returns $\boldsymbol{\mu}$ and
the covariance matrix $\boldsymbol{\Sigma}$. This work formulates the model as a
convex quadratic program (QP), solves it exactly through its Karush-Kuhn-Tucker
(KKT) conditions, and empirically assesses how far statistical and
\textit{Machine Learning} techniques improve the estimation of these parameters.
Using a sample of 30 S\&P~500 stocks (2010--2024), three covariance estimators
(sample, Ledoit-Wolf \textit{shrinkage} and principal-component (PCA)
reconstruction) and three models for predicting $\boldsymbol{\mu}$ (Ridge, Lasso
and Random Forest) are compared, all validated through walk-forward
\textit{backtesting} without \textit{look-ahead} against the $1/N$ (equally
weighted) and market-index (SPY) baselines. Shrinkage and PCA modestly improve
the sample portfolio (Sharpe $0.68$ and $0.69$ vs.\ $0.64$). When trained in a
pooled fashion, the \textit{Machine Learning} models are competitive: Lasso
(Sharpe $0.75$) matches the $1/N$ strategy ($0.76$). Nevertheless, no optimized
portfolio beats the market (Sharpe $0.90$) and, crucially, none of the Sharpe
differences is statistically significant. The main contribution is not a model
that beats the market, but a rigorous demonstration of why estimation error
dominates the problem, to the point of making the strategies indistinguishable,
and of which techniques help to mitigate it.

\keywordsen{Portfolio optimization, Markowitz model, Convex quadratic programming, Ledoit-Wolf shrinkage, Machine Learning, Walk-forward backtesting.}

\MediaOptionLogicBlank